\FloatBarrier
\section{نتیجه‌های \lr{v3}}
\label{sec:v3-results}

همه‌ی عددهای این فصل از جدول‌های ذخیره‌شده‌ی گام‌های ۲، ۳، ۵ و ۶ می‌آیند (\lr{results/v3} و \lr{outputs/checkpoints\_v3}). هر میانگین روی ۱۲ نفر است و $\pm$ خطای معیار میانگین را نشان می‌دهد.

\subsection{ساعت تکلیف بدون مغز}
\label{sec:v3-timing}

جدول~\ref{tab:v3-timing} و شکل~\ref{fig:v3_timing.png} نشان می‌دهند مدل‌هایی که فقط زمان رخدادها را می‌دانند تا کجا می‌رسند. مدل‌های \lr{cue} و \lr{phase} عددهای \lr{v2} را دقیقاً تکرار می‌کنند (\lr{۰٫۵۷۶} و \lr{۰٫۸۱۴}). نسخه‌های علّی کمی پایین‌ترند، چون رخدادهای آینده را نمی‌بینند: \lr{events\_causal} در افق صفر به \lr{۰٫۷۶۴} می‌رسد و \lr{causal\_cue} به \lr{۰٫۵۳۹}.

دو روند جالب است. دقت \lr{causal\_cue} با بلندتر شدن افق بالا می‌رود (از \lr{۰٫۵۳۹} به \lr{۰٫۶۱۹})، چون حرکت چند صد میلی‌ثانیه بعد از چراغ شروع می‌شود: زمان چراغ درباره‌ی فعالیت آینده بیشتر می‌گوید تا فعالیت همان لحظه. دقت \lr{events\_causal} برعکس کمی پایین می‌آید (به \lr{۰٫۷۲۹} در یک ثانیه)، چون رخدادهای بعدی هنوز اتفاق نیفتاده‌اند. در هر حال، یک مدل خطی بدون هیچ سیگنال مغزی یا حرکتی به $r$ حدود \lr{۰٫۷۵} می‌رسد. این بالاتر از \lr{E} مدل \lr{v2} است (\lr{۰٫۶۴۵}).

\begin{table}[htbp]
\centering\small
\caption{میانگین $r$ مدل‌های زمان‌بندی روی سری‌های آزمون، امتیاز در ۵۰۰ هرتز (عددهای ۵۰ هرتز تا سه رقم اعشار یکسان‌اند).}
\label{tab:v3-timing}
\begin{tabular}{lcccc}
\toprule
مدل & \lr{۰} میلی‌ثانیه & \lr{۲۰۰} & \lr{۵۰۰} & \lr{۱۰۰۰} \\
\midrule
\lr{cue} (غیرعلّی، سقف) & \lr{۰٫۵۷۶} $\pm$ \lr{۰٫۰۲۰} & -- & -- & -- \\
\lr{phase} (غیرعلّی، سقف) & \lr{۰٫۸۱۴} $\pm$ \lr{۰٫۰۱۵} & -- & -- & -- \\
\lr{causal\_cue} & \lr{۰٫۵۳۹} $\pm$ \lr{۰٫۰۲۳} & \lr{۰٫۵۵۸} $\pm$ \lr{۰٫۰۲۴} & \lr{۰٫۵۸۸} $\pm$ \lr{۰٫۰۲۴} & \lr{۰٫۶۱۹} $\pm$ \lr{۰٫۰۲۱} \\
\lr{events\_causal} & \lr{۰٫۷۶۴} $\pm$ \lr{۰٫۰۱۴} & \lr{۰٫۷۶۱} $\pm$ \lr{۰٫۰۱۴} & \lr{۰٫۷۵۳} $\pm$ \lr{۰٫۰۱۵} & \lr{۰٫۷۲۹} $\pm$ \lr{۰٫۰۱۶} \\
\bottomrule
\end{tabular}

\end{table}

\fig{v3_timing.png}{0.95}{مدل‌های زمان‌بندی بدون \lr{EEG}. چپ: میانگین $r$ در برابر افق. راست: هر عضله در افق صفر. ترتیب عضله‌ها همان ترتیب \lr{E} در \lr{v2} است (شانه بهتر از انگشت‌ها).}

\subsection{مقایسه‌ی اصلی: \lr{F} از \lr{B} بهتر نیست}
\label{sec:v3-bvsf}

جدول~\ref{tab:v3-bvsf} نتیجه‌ی اصلی را نشان می‌دهد. در هیچ‌کدام از چهار افق، \lr{F} از \lr{B} بهتر نیست. \lr{eeg\_skill} در همه‌ی افق‌ها کمی منفی است، یعنی افزودن \lr{EEG} خطای پیش‌بینی را کمی بیشتر می‌کند، بین \lr{۰٫۶} و \lr{۲٫۰} درصد. هیچ‌کدام از صفر فاصله‌ی معناداری ندارند: $p$ یک‌طرفه بین \lr{۰٫۷۲} و \lr{۰٫۸۸} است و پس از \lr{Holm} برابر \lr{۱٫۰۰}. در سطح هر نفر، فقط ۲ یا ۳ نفر از ۱۲ نفر \lr{eeg\_skill} مثبت و معنادار دارند، که از شانس انتظار می‌رود. \lr{VAF} و \lr{CCC} همین تصویر را تأیید می‌کنند: اختلاف \lr{VAF} بین \lr{-۰٫۰۹} و \lr{-۰٫۸۶} واحد است و \lr{CCC} دو مدل تا سه رقم اعشار تقریباً یکی است.

\begin{table}[htbp]
\centering\small
\caption{مقایسه‌ی اصلی روی همه‌ی \lr{token}های آزمون، با \lr{EEG} واقعی. «معنادار در» تعداد نفرهایی است که \lr{eeg\_skill} مثبت و $p$ کمتر از \lr{۰٫۰۵} دارند.}
\label{tab:v3-bvsf}
\resizebox{\textwidth}{!}{\begin{tabular}{lccccccc}
\toprule
افق (میلی‌ثانیه) & \lr{eeg\_skill} & $p$ & $p$ پس از \lr{Holm} & معنادار در & \lr{VAF gain} & \lr{CCC} مدل \lr{B} & \lr{CCC} مدل \lr{F} \\
\midrule
\lr{۰} & \lr{-۰٫۰۰۸۴} $\pm$ \lr{۰٫۰۰۷۱} & \lr{۰٫۸۳} & \lr{۱٫۰۰} & \lr{۲} & \lr{-۰٫۰۹} & \lr{۰٫۸۷۶} & \lr{۰٫۸۷۷} \\
\lr{۲۰۰} & \lr{-۰٫۰۱۴۵} $\pm$ \lr{۰٫۰۱۱۰} & \lr{۰٫۸۵} & \lr{۱٫۰۰} & \lr{۳} & \lr{-۰٫۱۹} & \lr{۰٫۸۵۶} & \lr{۰٫۸۵۴} \\
\lr{۵۰۰} & \lr{-۰٫۰۲۰۱} $\pm$ \lr{۰٫۰۱۶۲} & \lr{۰٫۸۸} & \lr{۱٫۰۰} & \lr{۲} & \lr{-۰٫۸۶} & \lr{۰٫۸۳۳} & \lr{۰٫۸۳۰} \\
\lr{۱۰۰۰} & \lr{-۰٫۰۰۵۶} $\pm$ \lr{۰٫۰۰۹۳} & \lr{۰٫۷۲} & \lr{۱٫۰۰} & \lr{۳} & \lr{-۰٫۲۵} & \lr{۰٫۷۸۴} & \lr{۰٫۷۸۰} \\
\bottomrule
\end{tabular}
}
\end{table}

\fig{v3_bvsf_horizon.png}{0.85}{میانگین \lr{eeg\_skill} در برابر افق برای هر منبع \lr{EEG}. خط پررنگ آبی \lr{EEG} واقعی است. هیچ نقطه‌ای ستاره (معناداری پس از \lr{Holm}) ندارد و همه‌ی منابع، از جمله \lr{EEG} جابه‌جاشده، در یک محدوده‌اند.}

\paragraph{کنترل‌ها همان رفتار را دارند.} جدول~\ref{tab:v3-sources} نشان می‌دهد \lr{EEG} واقعی از \lr{EEG} جابه‌جاشده بهتر نیست. فاصله‌ی \lr{eeg} با \lr{phase\_swap} در هیچ افقی معنادار نیست، و در افق ۲۰۰ میلی‌ثانیه \lr{EEG} جابه‌جاشده حتی کمی بهتر است. پس اطلاعاتی مخصوص هر \lr{trial} در \lr{EEG} دیده نمی‌شود. تنها عدد مثبت قابل‌توجه، توان باند کانال‌های پیرامونی در افق یک ثانیه است (\lr{+۰٫۰۲۰})، که آن هم معنادار نیست ($p$ خام \lr{۰٫۰۹} و پس از \lr{Holm} \lr{۰٫۳۵}) و اگر بود، به‌جای قشر حرکتی به \lr{artifact} اشاره می‌کرد.

\begin{table}[htbp]
\centering\small
\caption{میانگین \lr{eeg\_skill} برای هر منبع \lr{EEG} ورودی \lr{F}، همه‌ی \lr{token}های آزمون.}
\label{tab:v3-sources}
\begin{tabular}{lcccc}
\toprule
منبع \lr{EEG} در \lr{F} & \lr{۰} میلی‌ثانیه & \lr{۲۰۰} & \lr{۵۰۰} & \lr{۱۰۰۰} \\
\midrule
\lr{EEG} واقعی & \lr{-۰٫۰۰۸۴} & \lr{-۰٫۰۱۴۵} & \lr{-۰٫۰۲۰۱} & \lr{-۰٫۰۰۵۶} \\
\lr{phase\_swap} & \lr{-۰٫۰۰۶۸} & \lr{-۰٫۰۱۲۵} & \lr{-۰٫۰۲۴۲} & \lr{-۰٫۰۱۴۳} \\
\lr{swap} & \lr{-۰٫۰۰۷۷} & \lr{-۰٫۰۲۰۸} & \lr{-۰٫۰۱۲۸} & \lr{-۰٫۰۰۷۳} \\
\lr{bandpower} & \lr{-۰٫۰۰۳۴} & \lr{-۰٫۰۲۴۶} & \lr{+۰٫۰۰۲۸} & \lr{+۰٫۰۰۷۶} \\
\lr{peripheral} & \lr{-۰٫۰۰۳۷} & \lr{-۰٫۰۱۲۹} & \lr{-۰٫۰۰۰۲} & \lr{+۰٫۰۱۹۶} \\
\bottomrule
\end{tabular}

\end{table}

\paragraph{در هیچ فاز تکلیفی هم اثری نیست.} شکل~\ref{fig:v3_bvsf_phases.png} \lr{eeg\_skill} را برای هر فاز و هر افق نشان می‌دهد. هیچ خانه‌ای پس از \lr{Holm} معنادار نیست. اگر \lr{EEG} اطلاعات آمادگی حرکتی داشت، انتظار می‌رفت در فاز پیش از حرکت و افق‌های بلند مثبت شود. اما در همین فاز \lr{eeg\_skill} در هر چهار افق منفی است (بین \lr{-۰٫۰۱۴} و \lr{-۰٫۰۴۵}).

\fig{v3_bvsf_phases.png}{0.75}{\lr{eeg\_skill} مدل \lr{F} با \lr{EEG} واقعی برای هر فاز تکلیف و هر افق. ستاره یعنی معناداری پس از \lr{Holm} روی فازها؛ هیچ خانه‌ای ستاره ندارد.}

\paragraph{نتیجه به چند نفر خاص بسته نیست.} جدول~\ref{tab:v3-subjects} و شکل~\ref{fig:v3_bvsf_participants.png} مقدار هر نفر را نشان می‌دهند. در هر افق، \lr{eeg\_skill} فقط برای ۳ تا ۶ نفر مثبت است و هیچ نفری در همه‌ی افق‌ها مثبت نیست. بزرگ‌ترین مقدارهای منفی (مثل \lr{P6} در همه‌ی افق‌ها و \lr{P3} در ۵۰۰ میلی‌ثانیه) احتمالاً از بیش‌برازش \lr{ridge} روی ویژگی‌های \lr{EEG} می‌آیند، چون \lr{ridge} فقط روی یک سری \lr{validation} برازش می‌شود.

\begin{table}[htbp]
\centering\small
\caption{\lr{eeg\_skill} هر نفر با \lr{EEG} واقعی، همه‌ی \lr{token}ها.}
\label{tab:v3-subjects}
\begin{tabular}{lcccc}
\toprule
شرکت‌کننده & \lr{۰} میلی‌ثانیه & \lr{۲۰۰} & \lr{۵۰۰} & \lr{۱۰۰۰} \\
\midrule
\lr{P1} & \lr{-۰٫۰۱۲} & \lr{+۰٫۰۱۳} & \lr{-۰٫۰۲۹} & \lr{-۰٫۰۱۰} \\
\lr{P2} & \lr{-۰٫۰۰۳} & \lr{+۰٫۰۰۹} & \lr{+۰٫۰۴۴} & \lr{-۰٫۰۲۱} \\
\lr{P3} & \lr{+۰٫۰۰۶} & \lr{-۰٫۰۳۲} & \lr{-۰٫۱۴۳} & \lr{-۰٫۰۱۱} \\
\lr{P4} & \lr{-۰٫۰۱۰} & \lr{+۰٫۰۰۱} & \lr{-۰٫۰۱۵} & \lr{-۰٫۰۳۸} \\
\lr{P5} & \lr{+۰٫۰۰۳} & \lr{-۰٫۰۰۴} & \lr{-۰٫۰۱۱} & \lr{-۰٫۰۲۶} \\
\lr{P6} & \lr{-۰٫۰۶۳} & \lr{-۰٫۱۲۰} & \lr{-۰٫۰۶۹} & \lr{-۰٫۰۶۴} \\
\lr{P7} & \lr{-۰٫۰۳۷} & \lr{-۰٫۰۲۶} & \lr{+۰٫۰۷۰} & \lr{+۰٫۰۵۰} \\
\lr{P8} & \lr{+۰٫۰۲۶} & \lr{+۰٫۰۰۶} & \lr{-۰٫۰۴۵} & \lr{+۰٫۰۳۳} \\
\lr{P9} & \lr{-۰٫۰۲۸} & \lr{-۰٫۰۳۲} & \lr{+۰٫۰۳۱} & \lr{+۰٫۰۱۴} \\
\lr{P10} & \lr{+۰٫۰۰۸} & \lr{+۰٫۰۱۹} & \lr{-۰٫۰۵۶} & \lr{+۰٫۰۲۵} \\
\lr{P11} & \lr{-۰٫۰۰۶} & \lr{+۰٫۰۱۴} & \lr{-۰٫۰۱۲} & \lr{+۰٫۰۰۰} \\
\lr{P12} & \lr{+۰٫۰۱۶} & \lr{-۰٫۰۲۲} & \lr{-۰٫۰۰۴} & \lr{-۰٫۰۲۰} \\
\midrule
مثبت در & \lr{۵} از \lr{۱۲} & \lr{۶} از \lr{۱۲} & \lr{۳} از \lr{۱۲} & \lr{۵} از \lr{۱۲} \\
\bottomrule
\end{tabular}

\end{table}

\fig{v3_bvsf_participants.png}{0.8}{\lr{eeg\_skill} هر نفر (خاکستری) و میانگین (آبی) در برابر افق.}

\subsection{هر گام از ساعت تکلیف تا \lr{EEG} چقدر خطا را کم می‌کند}

جدول~\ref{tab:v3-mse} خطای مطلق هر شرط را کنار هم می‌گذارد و شکل~\ref{fig:v3_mse_ladder.png} همین را نشان می‌دهد. سه نکته از آن پیداست:

\begin{itemize}
\item سینماتیک بیشترین سهم را دارد. خطای \lr{K} در افق صفر حدود یک‌پنجم خطای \lr{causal\_cue} است. با بلندتر شدن افق، خطای \lr{K} بیشتر می‌شود (از \lr{۰٫۰۰۶۹} به \lr{۰٫۰۱۶۵})، چون حرکت گذشته کمتر درباره‌ی آینده‌ی دورتر می‌گوید.
\item رخدادهای گذشته فقط در افق‌های میانی کمک می‌کنند. در ۲۰۰ و ۵۰۰ میلی‌ثانیه، \lr{B} خطای \lr{K} را حدود ۳ و ۲ درصد کم می‌کند. در افق صفر \lr{B} حدود ۲ درصد بدتر از \lr{K} است و در یک ثانیه تفاوتی ندارد. پس وقتی رخدادها چیزی اضافه نمی‌کنند، \lr{ridge} هم کمی بیش‌برازش می‌کند.
\item افزودن \lr{EEG} خطا را در هر افق کمی بیشتر می‌کند (\lr{+۰٫۷} تا \lr{+۲٫۵} درصد نسبت به \lr{B}).
\end{itemize}

\begin{table}[htbp]
\centering\small
\caption{میانگین خطای مطلق (\lr{MSE} در حوزه‌ی مدل، میانگین پنج عضله) هر شرط روی همه‌ی \lr{token}های آزمون.}
\label{tab:v3-mse}
\resizebox{\textwidth}{!}{\begin{tabular}{lcccccc}
\toprule
افق (میلی‌ثانیه) & \lr{causal\_cue} & \lr{K} & \lr{B} & \lr{F} & تغییر \lr{B} نسبت به \lr{K} (درصد) & تغییر \lr{F} نسبت به \lr{B} (درصد) \\
\midrule
\lr{۰} & \lr{۰٫۰۳۵۹} & \lr{۰٫۰۰۶۹۲} & \lr{۰٫۰۰۷۰۶} & \lr{۰٫۰۰۷۱۲} & \lr{+۲٫۱} & \lr{+۰٫۸} \\
\lr{۲۰۰} & \lr{۰٫۰۳۳۶} & \lr{۰٫۰۰۸۶۴} & \lr{۰٫۰۰۸۳۹} & \lr{۰٫۰۰۸۵۱} & \lr{-۲٫۸} & \lr{+۱٫۴} \\
\lr{۵۰۰} & \lr{۰٫۰۳۰۶} & \lr{۰٫۰۱۱۱۴} & \lr{۰٫۰۱۰۸۸} & \lr{۰٫۰۱۱۱۵} & \lr{-۲٫۳} & \lr{+۲٫۵} \\
\lr{۱۰۰۰} & \lr{۰٫۰۲۸۱} & \lr{۰٫۰۱۶۴۶} & \lr{۰٫۰۱۶۵۰} & \lr{۰٫۰۱۶۶۱} & \lr{+۰٫۲} & \lr{+۰٫۷} \\
\bottomrule
\end{tabular}
}
\end{table}

\fig{v3_mse_ladder.png}{0.8}{خطای مطلق \lr{causal\_cue}، \lr{K}، \lr{B} و \lr{F} در هر افق.}

\subsection{زمان شروع فعالیت عضله}

جدول~\ref{tab:v3-onset} میانه‌ی قدرمطلق خطای \lr{onset} را نشان می‌دهد. زمان‌ها در نرخ \lr{token} (۵۰ هرتز) سنجیده شده‌اند، پس دقت آن‌ها ۲۰ میلی‌ثانیه است. سینماتیک زمان شروع را خیلی بهتر از چراغ تنها پیش‌بینی می‌کند (۸۰ در برابر ۳۲۰ میلی‌ثانیه در افق صفر). در افق‌های ۲۰۰ و ۵۰۰ میلی‌ثانیه، رخدادهای گذشته خطای \lr{K} را کم می‌کنند (در ۵۰۰ میلی‌ثانیه از ۴۰۰ به ۲۴۰ میلی‌ثانیه). \lr{F} در هیچ افقی زمان شروع را بهتر از \lr{B} پیش‌بینی نمی‌کند.

\begin{table}[htbp]
\centering\small
\caption{میانه‌ی قدرمطلق خطای زمان شروع فعالیت (میلی‌ثانیه)، روی همه‌ی \lr{lift}ها و عضله‌ها.}
\label{tab:v3-onset}
\begin{tabular}{lcccc}
\toprule
افق (میلی‌ثانیه) & \lr{causal\_cue} & \lr{K} & \lr{B} & \lr{F} \\
\midrule
\lr{۰} & \lr{۳۲۰} & \lr{۸۰} & \lr{۸۰} & \lr{۸۰} \\
\lr{۲۰۰} & \lr{۲۸۰} & \lr{۲۲۰} & \lr{۱۸۰} & \lr{۲۰۰} \\
\lr{۵۰۰} & \lr{۲۸۰} & \lr{۴۰۰} & \lr{۲۴۰} & \lr{۲۶۰} \\
\lr{۱۰۰۰} & \lr{۵۸۰} & \lr{۴۲۰} & \lr{۴۰۰} & \lr{۴۰۰} \\
\bottomrule
\end{tabular}

\end{table}

\fig{v3_onset.png}{0.8}{میانه‌ی قدرمطلق خطای \lr{onset} هر شرط در برابر افق (میانگین روی نفرها).}

\subsection{آموزش \lr{E0} و \lr{Eguided}}
\label{sec:v3-training}

هر ۱۴۴ \lr{run} آموزش تمام شده‌اند: \lr{E0} و سه مقدار \lr{scale}، هر کدام با سه \lr{seed} برای ۱۲ نفر. هر \lr{run} روی \lr{GTX 1660 Ti} به‌طور میانگین حدود ۱۳٫۵ دقیقه طول کشید و کل آموزش حدود ۳۲ ساعت \lr{GPU} بود. شکل~\ref{fig:v3_training.png} منحنی‌ها را نشان می‌دهد: وزن راهنما همان برنامه‌ی تعریف‌شده را دنبال می‌کند و سهم راهنما از گرادیان با \lr{scale} بزرگ‌تر می‌شود (بیشینه‌ی میانگین \lr{۰٫۰۷}، \lr{۰٫۱۴} و \lr{۰٫۳۴}). پس راهنما واقعاً روی آموزش اثر گذاشته است.

اما این اثر به \lr{loss} \lr{validation} نمی‌رسد (جدول~\ref{tab:v3-training}). بهترین \lr{loss} \lr{validation} مدل‌های راهنماشده با \lr{E0} عملاً برابر است. بیشترین تغییر با \lr{scale} \lr{۰٫۳} است (\lr{-۰٫۸۳} درصد)، که در ۷ نفر از ۱۲ کمتر از \lr{E0} است و معنادار نیست ($p$ برابر \lr{۰٫۱۰}).

\begin{table}[htbp]
\centering\small
\caption{آموزش \lr{E0} و \lr{Eguided}: میانگین روی ۳۶ \lr{run} هر مدل. تغییر \lr{loss} و مقایسه با \lr{E0} روی میانگین سه \lr{seed} هر نفر است؛ $p$ از آزمون \lr{Wilcoxon} یک‌طرفه (کمتر بودن \lr{loss}).}
\label{tab:v3-training}
\resizebox{\textwidth}{!}{\begin{tabular}{lcccccccc}
\toprule
مدل & \lr{run} & \lr{epoch} & بهترین \lr{epoch} & \lr{loss} \lr{validation} & بیشینه‌ی سهم گرادیان راهنما & تغییر \lr{loss} (درصد) & کمتر از \lr{E0} & $p$ \\
\midrule
\lr{E0} & \lr{۳۶} & \lr{۸۳} & \lr{۵۳} & \lr{۰٫۱۹۹۰} & \lr{۰٫۰۰} & -- & -- & -- \\
\lr{Eguided}، \lr{scale} \lr{۰٫۰۵} & \lr{۳۶} & \lr{۸۶} & \lr{۵۵} & \lr{۰٫۱۹۹۰} & \lr{۰٫۰۷} & \lr{+۰٫۰۴} & \lr{۷} از \lr{۱۲} & \lr{۰٫۴۰} \\
\lr{Eguided}، \lr{scale} \lr{۰٫۱} & \lr{۳۶} & \lr{۸۷} & \lr{۵۵} & \lr{۰٫۱۹۸۷} & \lr{۰٫۱۴} & \lr{-۰٫۱۳} & \lr{۷} از \lr{۱۲} & \lr{۰٫۵۲} \\
\lr{Eguided}، \lr{scale} \lr{۰٫۳} & \lr{۳۶} & \lr{۸۱} & \lr{۴۹} & \lr{۰٫۱۹۷۳} & \lr{۰٫۳۴} & \lr{-۰٫۸۳} & \lr{۷} از \lr{۱۲} & \lr{۰٫۱۰} \\
\bottomrule
\end{tabular}
}
\end{table}

\fig{v3_training.png}{0.98}{منحنی‌های آموزش هر ۱۴۴ \lr{run}. چپ: \lr{loss} \lr{validation} خود \lr{E}، که انتخاب مدل فقط با آن انجام می‌شود. وسط: وزن راهنما. راست: سهم راهنما از گرادیان.}

\subsection{آنچه هنوز در راه است}
\label{sec:v3-pending}

گام ۷ (ارزیابی \lr{E0} و \lr{Eguided} روی سری‌های آزمون) هنگام نوشتن این گزارش در حال اجراست. این گام برای هر \lr{run} این‌ها را می‌سنجد: دقت در ۵۰۰ و ۵۰ هرتز، آزمون‌های \lr{null}، \lr{phase-matched swap}، زمان شروع فعالیت، باقی‌مانده نسبت به هر مدل زمان‌بندی، و آزمون انحراف از الگو. پس از آن گام ۸ حکم نهایی را با هر دو پرسش می‌دهد. نتیجه‌ی \lr{validation} نشان می‌دهد که احتمالاً \lr{Eguided} از \lr{E0} به‌طور معنادار بهتر نخواهد شد. اگر چنین باشد، با توجه به نتیجه‌ی مقایسه‌ی اصلی، حکم نهایی سطر آخر جدول~\ref{tab:v3-interp} خواهد بود: شاهدی نیست که \lr{EEG} چیزی فراتر از حرکت و زمان‌بندی داشته باشد. این یک پیش‌بینی است، نه نتیجه.

\subsection{بحث}

\paragraph{نتیجه‌ی اصلی.} وقتی حرکت گذشته و زمان رخدادهای گذشته را داریم، تاریخچه‌ی \lr{EEG} در هیچ افقی از صفر تا یک ثانیه، و در هیچ فاز تکلیفی، چیزی به پیش‌بینی پوش \lr{EMG} اضافه نمی‌کند. کنترل‌ها همین را تأیید می‌کنند: \lr{EEG} یک \lr{trial} دیگر با همان زمان‌بندی همان‌قدر خوب است.

\paragraph{ارتباط با \lr{v2}.} در \lr{v2} نسخه‌ی علّی، \lr{F} در افق ۵۰۰ میلی‌ثانیه حدود \lr{۰٫۵۵} واحد \lr{VAF} از \lr{K} بهتر بود. در \lr{v3}، مبنا رخدادهای گذشته را هم دارد و این برتری از بین می‌رود. این با این فرض سازگار است که سهم \lr{EEG} در \lr{v2} از دانستن زمان تکلیف می‌آمد، مثلاً از پاسخ \lr{EEG} به چراغ \lr{LED}، نه از فرمان حرکتی. این یک استنباط است: \lr{v3} نشان می‌دهد \lr{EEG} به \lr{B} چیزی اضافه نمی‌کند، ولی مستقیماً نمی‌سنجد سهم \lr{EEG} در \lr{v2} کدام بخش از اطلاعات رخدادها بوده است.

\paragraph{محدودیت‌ها.}
\begin{itemize}
\item \lr{ridge} مدل‌های \lr{B} و \lr{F} فقط روی یک سری \lr{validation} برازش می‌شود، چون \lr{K} روی سری‌های آموزش آموزش دیده است. افزودن ده‌ها ویژگی \lr{EEG} با این داده‌ی کم می‌تواند به بیش‌برازش و \lr{eeg\_skill} منفی برسد. یک اثر کوچک واقعی ممکن است زیر این هزینه پنهان بماند.
\item ویژگی‌های \lr{EEG} در \lr{F}، \lr{token}های جریان \lr{EEG} مدل علّی \lr{v2} هستند که برای پیش‌بینی خود \lr{EMG} آموزش دیده‌اند، نه برای پیش‌بینی خطای \lr{B}. این ریسک هست که اطلاعات مفیدی در این \lr{token}ها فشرده نشده باشد. منبع \lr{bandpower} که مستقل از شبکه است همان نتیجه را می‌دهد، ولی آن هم فقط خطی است.
\item رخدادهای تکلیف جز \lr{LED} از ضبط سینماتیک و نیرو استخراج شده‌اند. ویژگی‌ها فقط رخدادهای گذشته را می‌خوانند، ولی زمان آن‌ها از کل ضبط تشخیص داده شده است. \lr{gate} نیرو در مدل علّی هم حدود ۵۰ میلی‌ثانیه آینده را می‌بیند. هر دو به نفع \lr{B} هستند، پس نتیجه‌گیری درباره‌ی \lr{EEG} را محافظه‌کارانه‌تر می‌کنند.
\item با ۱۲ نفر، آزمون گروهی فقط اثرهای متوسط را به‌طور مطمئن می‌بیند.
\end{itemize}

\subsection{پیشنهاد برای جلسه}

\begin{enumerate}
\item پذیرفتن نتیجه‌ی اصلی به‌عنوان یافته‌ی منفی مقاله، همراه با نکته‌ی روش‌شناختی: دقت بالای مدل‌های \lr{EEG} به \lr{EMG} در این داده را می‌شود تقریباً کامل با حرکت و زمان‌بندی تکلیف توضیح داد، و هر ادعای سهم \lr{EEG} باید با مبنای \lr{B} و کنترل \lr{phase-matched swap} سنجیده شود.
\item یک آزمون حساسیت ارزان برای کنار گذاشتن محدودیت اول: برازش \lr{ridge} با \lr{cross-validation} روی \lr{trial}های سری‌های آزمون به‌جای یک سری \lr{validation}، تا داده‌ی برازش بیشتر شود.
\item تصمیم درباره‌ی کنترل‌های \lr{artifact} با آموزش (\lr{central}، \lr{peripheral} و \lr{eye\_regressed}). با نبود اثر مثبت، این‌ها اولویت کمتری دارند.
\item منتظر ماندن برای گام‌های ۷ و ۸ پیش از نهایی کردن پرسش ثانویه.
\end{enumerate}
